\par\Needspace{7\baselineskip}\SourceLeaf{405}
\begin{center}\textsc{Gauss} an \textsc{Olbers}. Göttingen, 30. Oktober 1825.\end{center}
\indent\mbox{. . . . . . . . . .} Über meine Bedenken zur Gliederung meines noch auszuarbeitenden Werks über Höhere Geodäsie habe ich mich in meinem Brief wohl nicht klar ausgedrückt. Es geht dabei nicht um die Frage, ob ich meine Auffassung von den Krümmungsradien aufnehmen soll. Es geht um die immer umfangreicher werdenden ganz allgemeinen Untersuchungen über die krummen Flächen, die auf ihnen gebildeten Figuren, die Natur der kürzesten und nicht kürzesten Linien und eine Menge anderer Gegenstände, die ich hier nicht anführen kann, weil die Begriffe noch nicht gebräuchlich sind und es nicht einmal Namen dafür gibt --- ob ich all dies aufnehmen soll, zumal ich täglich mehr Themen finde, auf die ich ebenfalls zurückgehen muss. Meines Wissens wurden sie bisher nicht unter dem Gesichtspunkt betrachtet, der für den Zusammenhang des Ganzen nötig ist. Dazu gehörte (als ein unbedeutendes Beispiel) sogar die Lehre von den Krümmungsradien in Plano.

Ich hätte auch die Lehre vom Flächeninhalt der Figuren im Allgemeinen nennen können, die ich ebenfalls seit 30 und mehr Jahren unter einem Gesichtspunkt betrachtet habe, den ich bisher für neu hielt. Das war zum Teil ein Irrtum. Erst vor Kurzem habe ich eine Abhandlung von \textsc{Meister} (meiner Meinung nach ein sehr genialer Kopf) im ersten Band der Novi Commentarii Gotting. kennengelernt, in der die Sache fast genauso betrachtet und sehr schön entwickelt wird.

Diese treffliche Abhandlung ist den Mathematikern aber fast ganz unbekannt. Es würde auch nicht ausreichen, mich auf sie zu beziehen, da sie nur die ersten Grundbegriffe enthält. Bei der analytischen Behandlung geometrischer Beziehungen hat man zwar längst Linien von positivem und negativem Wert gut verstanden und erkannt, dass dabei immer ausdrücklich oder stillschweigend ein gewisser Sinn (sens, Richtung) zugrunde liegt, in dem die Linien als wachsend angesehen werden usw. Wenn man aber Flächeninhalte (areas) durch Formeln ausdrückt, muss natürlich auch ein %
